% !TEX program = xelatex
\documentclass[11pt, letterpaper]{article}

\usepackage{fontspec}
\setmainfont{Libre Caslon Text}

\usepackage[top=1in, bottom=1in, left=1in, right=1in]{geometry}
\usepackage{parskip}
\setlength{\parindent}{0pt}
\usepackage{amsmath}
\usepackage{booktabs}
\usepackage{xcolor}
\usepackage[hidelinks]{hyperref}

\definecolor{nycblue}{HTML}{003087}

\begin{document}

\begin{center}
  {\large\bfseries\color{nycblue} Superintendent's Portal Overall Match Rate Calculation Correction}\\[0.4em]
\end{center}

\bigskip

\textbf{Background.} The dashboard compares SubCentral job records against
SREPP payroll records on a per-location basis to produce a \emph{match rate}
--- the share of payroll days that have a corresponding SubCentral entry.
A discrepancy was identified between the headline ``Overall Match Rate'' and
the per-borough rates shown in the summary table.

\bigskip

\textbf{Previous formula (unweighted mean).}
The overall rate was previously computed as the arithmetic mean of each
school's individual match percentage:

\[
  \bar{r}_{\text{unweighted}}
  = \frac{1}{n} \sum_{i=1}^{n} \frac{m_i}{p_i} \times 100
\]

where $n$ is the number of school locations, $m_i$ is the number of matched
job days at location $i$, and $p_i$ is the number of payroll job days at
location $i$. Under this formula every school receives equal weight regardless
of size. Schools with few or no payroll records (including locations present
in SubCentral but absent from the payroll file, which contribute $0\%$) pull
the average downward disproportionately, yielding a headline figure several
percentage points below the borough-level rates.

\bigskip

\textbf{Current formula (weighted aggregate).}
The overall rate is now computed as a single ratio of citywide totals,
consistent with how per-borough rates are calculated:

\[
  r_{\text{aggregate}}
  = \frac{\displaystyle\sum_{i=1}^{n} m_i}{\displaystyle\sum_{i=1}^{n} p_i}
    \times 100
\]

Each payroll day contributes equally to the denominator, so larger schools
with more activity are weighted in proportion to their size. This is the same
methodology used for the per-borough rows in the summary table, making the
headline figure directly comparable to and consistent with those values.

\bigskip

\textbf{Illustrative example.}

\begin{center}
\renewcommand{\arraystretch}{1.3}
\begin{tabular}{lrrr}
\toprule
\textbf{School} & \textbf{Payroll Days} & \textbf{Matched Days} & \textbf{Per-School Rate} \\
\midrule
School A & 5{,}000 & 4{,}600 & 92.0\% \\
School B &    20   &     0   &  0.0\% \\
School C & 3{,}000 & 2{,}700 & 90.0\% \\
\midrule
\textbf{Unweighted mean}  & ---      & ---      & 60.7\% \\
\textbf{Weighted aggregate} & 8{,}020 & 7{,}300 & 91.0\% \\
\bottomrule
\end{tabular}
\end{center}

School B has minimal activity but drags the unweighted mean to 60.7\%,
far below the 91.0\% weighted aggregate that accurately reflects system-wide
coverage.

\bigskip

\textbf{Summary.} The per-borough rates displayed in the summary table have
always used the weighted aggregate formula and remain unchanged. Only the
headline ``Overall Match Rate'' figure has been corrected to use the same
methodology, resolving the apparent inconsistency between the headline and
the borough rows.

\bigskip

\textbf{Bottom line.} The previous calculation was averaging the match
percentage of each individual school --- including schools reporting $0\%$
--- giving every school equal weight regardless of how many payroll days it
contributed. The correct approach is to first add up all matched job days
across every school, then divide by the total payroll days citywide. This
produces a single percentage that reflects actual system-wide coverage rather
than a mean of school-level percentages distorted by low-activity locations.

\end{document}
